% !TeX root = ../../Book3.tex
\section{零点计数}
\label{b3:sec:1801}

辐角原理用一条闭曲线计算零点总数，但没有直接显示零点到圆心的距离。对数模的圆周平均会把这项几何信息补进计数。

\subsection{重数}
\label{b3:subsec:180101}

非恒零全纯函数在零点 \(a\) 可唯一写成 \(f(z)=(z-a)^m g(z)\)，其中 \(g(a)\ne0\)，\(m\) 是零点重数。后面的计数均按重数重复零点。紧子集上只有有限个零点；否则有内部聚点，恒等定理迫使函数恒零。

例如 \(f(z)=z^2(z-1)^3\) 在原点和一的质量分别为二、三。零点位置相同而重数不同，会改变对数模的奇性强度；按不同零点个数而不计重数不能恢复解析信息。

\subsection{计数函数}
\label{b3:subsec:180102}

对 \(f(0)\ne0\)，令 \(n(r)\) 为 \(|a|\le r\) 内按重数的零点数，定义积分计数
\[
N(r)=\int_0^r\frac{n(t)}t\,\mathrm{d}t
=\sum_{|a|<r}\log\frac r{|a|}.
\]
第二式由有限和换序得到。零点在半径边界时贡献零，故开闭圆盘约定不影响 \(N\)。若原点重数为 \(m\)，先除去 \(z^m\)，再单独保留 \(m\log r\) 项。\(n\) 是阶梯函数，\(N\) 却连续；每个零点从它进入圆盘时开始累计对数权重。

\subsection{辐角原理再述}
\label{b3:subsec:180103}

边界无零点时，\cref{b3:thm:03-argument-principle}给
\[
n(r)=\frac1{2\pi\mathrm{i}}\int_{|z|=r}\frac{f'(z)}{f(z)}\,\mathrm{d}z.
\]
这里要求 \(f\) 在闭圆盘邻域全纯，或者先在稍小半径应用；边界上有零点时积分不定义。辐角原理显示整数计数的拓扑稳定性，积分计数则更适合平均与增长估计。后面不会把 \(n(r)\) 的跳跃误当作对数平均本身的跳跃。


\ProseParagraphBreak

计数函数的两个版本表达不同信息。若半径内有 \(N\) 个零点全位于 \(|a|=r/2\)，它们给 \(n(r)=N\)、\(N(r)=N\log2\)；若相同个数全位于 \(|a|=(1-\varepsilon)r\)，整数计数仍为 \(N\)，积分计数却为 \(-N\log(1-\varepsilon)\)，随 \(\varepsilon\downarrow0\) 趋零。紧贴观察圆周的零点对该半径的平均影响较小，但一旦把半径扩大，它们会继续累积权重。因此从 \(N(r)\) 控制 \(n(r)\) 必须另取更大的半径，不能同半径直接比较而没有额外系数。

\ProseParagraphBreak

换一个圆心可以避开原点零点。对内部点 \(b\) 且 \(f(b)\ne0\)，在 \(B(b,R)\) 使用变量 \(\zeta=z-b\)，得到完全相同的计数和公式。如果要在几个不同圆心比较，应分别保留 \(\log|f(b)|\)；这个中心常数可能很负，会使零点上界较弱。例如函数在 \(b\) 非零但已非常接近一个零点时，\(|f(b)|\) 小，计数估计仍成立却不一定有效。归一化 \(f(b)=1\) 可消除常数项，但乘法归一化并不改变零点本身。

\ProseParagraphBreak

对含原点重数的函数，用修正积分计数
\[
N_f(r)=m\log r+\sum_{0<|a|<r}m_a\log(r/|a|)
\]
十分方便。它在小于一的半径可能为负，因为 \(m\log r<0\)；这不是非原点零点贡献出现负值，而是已单独记录的中心奇性。若只需要非负零点数的估计，先除去原点幂再应用公式更稳妥。所有符号版本必须在中心奇性项上保持同一约定，尤其进入亚纯函数的零点、极点差计数时不能重复加上原点项。

\ProseParagraphBreak

零点对全纯扰动的稳定性有一个比总数更精细的表述。设 \(f_j\to f\) 在区域 \(D\) 内局部一致，且 \(f\not\equiv0\)。固定一个内部零点 \(a\)，取小圆盘使闭盘中没有其他极限零点，边界上的 \(|f|\) 有正下界。Rouch\'e 定理说明：充分大的 \(j\) 在盘内恰有 \(m_a\) 个零点，按重数计。把圆盘继续缩小，仍有同样结论，所以这些零点全向 \(a\) 集中。紧集上远离极限零点的部分，由 \(|f|\) 的正下界可知最终没有新的零点。
\ProseParagraphBreak
于是离散测度 \(\mu_j=\sum m_{j,a}\delta_a\) 在内部弱收敛到 \(\mu=\sum m_a\delta_a\)：对紧支集连续测试函数，只需用有限个互不相交小圆盘包住相关极限零点，每盘总质量相同，测试函数在盘内的振幅任意小，其余部分最终无质量。这里弱收敛只针对内部紧支集；零点可以逃向边界或无穷远。例如 \(f_j(z)=1-z/j\) 在全平面局部一致趋于一，唯一零点趋向无穷远。极限函数没有零点，与每个近似函数有一个零点并不矛盾。
\ProseParagraphBreak
假设极限不恒零是必要的。\(f_j(z)=e^{-j}\sin(jz)\) 在足够窄的固定水平带中局部一致趋于零，但零点间距趋于零，任何固定实区间内的零点数都趋向无穷。恒等为零的函数不再拥有离散的零点测度，不能给它随意指定一个有限重数来延续上面的结论。
